\DeclareRobustCommand{\code}[1]{\nolinkurl{#1}}
\DeclareRobustCommand{\term}[2]{\textbf{#1}\index{#1}\,（#2）}
\DeclareRobustCommand{\engterm}[1]{\textbf{#1}\index{#1}}
\DeclareRobustCommand{\filepath}[1]{\nolinkurl{#1}}
\DeclareRobustCommand{\cmd}[1]{\nolinkurl{#1}}
\newcommand{\register}[1]{\texttt{#1}\index{register!#1}}
\newcommand{\instruction}[1]{\texttt{#1}\index{instruction!#1}}
\newcommand{\concept}[1]{\textbf{#1}\index{#1}}

\pdfstringdefDisableCommands{%
  \def\code#1{\detokenize{#1}}%
  \def\term#1#2{#1}%
  \def\engterm#1{#1}%
  \def\filepath#1{#1}%
  \def\cmd#1{#1}%
  \def\register#1{#1}%
  \def\instruction#1{#1}%
  \def\concept#1{#1}%
}

\newcommand{\plannedchapter}[5]{
  \section{#1}

  \subsection{本节定位}
  #2

  \subsection{核心问题}
  \begin{itemize}
    #3
  \end{itemize}

  \subsection{必须讲清楚的底层机制}
  \begin{itemize}
    #4
  \end{itemize}

  \subsection{实验和交付物}
  \begin{labbox}
  #5
  \end{labbox}

  \subsection{扩写标准}
  本节后续扩写时必须从具体 C/C++ 例子出发，逐层连接到编译器表示、汇编、机器执行、系统机制和性能证据。不能只列概念；每个核心概念都必须解释它为什么存在、没有它会发生什么、底层如何实现、如何用工具观察、对性能有什么影响。
}

\newcommand{\bookterm}[1]{\textbf{\color{BookNavy}#1}}
\newcommand{\important}[1]{\textbf{\color{BookRed}#1}}
\newcolumntype{L}[1]{>{\raggedright\arraybackslash}p{#1}}
